\usepackage[a4paper,width=15cm,height=24cm]{geometry}

\usepackage{amsmath,amsthm,amssymb}
\usepackage{bbm}
\usepackage{mathrsfs}

%\usepackage[authoryear,longnamesfirst]{natbib}
%\usepackage{array}
\usepackage{natbib}
\usepackage{color}
\usepackage[dvipsnames]{xcolor}
\usepackage{tikz}
\usepackage{tikz-cd}
\usetikzlibrary{shapes.geometric}
\usepackage{subcaption}
\usepackage{caption}


\usepackage{filemod}
\usepackage[breaklinks=true]{hyperref}

%\usepackage{showlabels}

% define theorem environments
\theoremstyle{plain}
\newtheorem{theorem}{Theorem}[section]
\newtheorem{proposition}[theorem]{Proposition}
\newtheorem{lemma}[theorem]{Lemma}
\newtheorem{corollary}[theorem]{Corollary}
\newtheorem{conjecture}[theorem]{Conjecture}

\theoremstyle{definition}
\newtheorem{question}[theorem]{Question}
\newtheorem{definition}[theorem]{Definition}
\newtheorem{example}[theorem]{Example}
\newtheorem{conjexample}[theorem]{Conjectural Example}
\newtheorem{remark}[theorem]{Remark}
\newtheorem{case}[theorem]{Case}
\newtheorem{condition}[theorem]{Condition}
\newtheorem{notation}[theorem]{Notation}


% Reset proofclaim counter at start of each proof
\newtheorem{proofclaim}{Claim}
\let\oldproof\proof
\renewcommand{\proof}{%
    \setcounter{proofclaim}{0}%
    \oldproof%
}


\makeatletter

%% display environments
\def\be#1{\begin{equation*}#1\end{equation*}}
\def\ben#1{\begin{equation}#1\end{equation}}
\def\bes#1{\begin{equation*}\begin{split}#1\end{split}\end{equation*}}
\def\besn#1{\begin{equation}\begin{split}#1\end{split}\end{equation}}
\def\bg#1{\begin{gather*}#1\end{gather*}}
\def\bgn#1{\begin{gather}#1\end{gather}}
\def\bm#1{\begin{multline*}#1\end{multline*}}
\def\bmn#1{\begin{multline}#1\end{multline}}
\def\ba#1{\begin{align*}#1\end{align*}}
\def\ban#1{\begin{align}#1\end{align}}

%%% bracket commands
\def\given{\typeout{Command 'given' should only be used within bracket command}}
\newcounter{@bracketlevel}
\def\@bracketfactory#1#2#3#4#5#6{
\expandafter\def\csname#1\endcsname##1{%
\addtocounter{@bracketlevel}{1}%
\global\expandafter\let\csname @middummy\alph{@bracketlevel}\endcsname\given%
\global\def\given{\mskip#5\csname#4\endcsname\vert\mskip#6}\csname#4l\endcsname#2##1\csname#4r\endcsname#3%
\global\expandafter\let\expandafter\given\csname @middummy\alph{@bracketlevel}\endcsname
\addtocounter{@bracketlevel}{-1}}%
}
\def\bracketfactory#1#2#3{%
\@bracketfactory{#1}{#2}{#3}{relax}{1mu plus 0.25mu minus 0.25mu}{0.6mu plus 0.15mu minus 0.15mu}
\@bracketfactory{b#1}{#2}{#3}{big}{1mu plus 0.25mu minus 0.25mu}{0.6mu plus 0.15mu minus 0.15mu}
\@bracketfactory{bb#1}{#2}{#3}{Big}{2.4mu plus 0.8mu minus 0.8mu}{1.8mu plus 0.6mu minus 0.6mu}
\@bracketfactory{bbb#1}{#2}{#3}{bigg}{3.2mu plus 1mu minus 1mu}{2.4mu plus 0.75mu minus 0.75mu}
\@bracketfactory{bbbb#1}{#2}{#3}{Bigg}{4mu plus 1mu minus 1mu}{3mu plus 0.75mu minus 0.75mu}
}
\bracketfactory{clc}{\lbrace}{\rbrace}
\bracketfactory{clr}{(}{)}
\bracketfactory{cls}{[}{]}
\bracketfactory{abs}{\lvert}{\rvert}
\bracketfactory{norm}{\Vert}{\Vert}
\bracketfactory{floor}{\lfloor}{\rfloor}
\bracketfactory{ceil}{\lceil}{\rceil}
\bracketfactory{angle}{\langle}{\rangle}

%% define \sA \cA \bA \IA for all capital letters A ... Z %%%%%%%
\newcounter{ctr}\loop\stepcounter{ctr}\edef\X{\@Alph\c@ctr}%
	\expandafter\edef\csname s\X\endcsname{\noexpand\mathscr{\X}}
	\expandafter\edef\csname c\X\endcsname{\noexpand\mathcal{\X}}
	\expandafter\edef\csname b\X\endcsname{\noexpand\boldsymbol{\X}}
	\expandafter\edef\csname I\X\endcsname{\noexpand\mathbbm{\X}}
	\expandafter\edef\csname r\X\endcsname{\noexpand\mathrm{\X}}
\ifnum\thectr<26\repeat

%%%%%%%%%%%%%%%%%%%%%%%%%%%%%%%

\newcount\minute
\newcount\hour
\newcount\hourMins
\def\now{%
\minute=\time%
\hour=\time \divide \hour by 60%
\hourMins=\hour \multiply\hourMins by 60%
\advance\minute by -\hourMins%
\zeroPadTwo{\the\hour}:\zeroPadTwo{\the\minute}%
}
\def\zeroPadTwo#1{\ifnum #1<10 0\fi#1}

% Change layout
\numberwithin{equation}{section}
\allowdisplaybreaks[4]

\renewcommand\section{\@startsection {section}{1}{\z@}%
{-3.5ex \@plus -1ex \@minus -.2ex}%
{1.3ex \@plus.2ex}%
{\center\small\sc\mathversion{bold}\MakeUppercase}}

\def\subsection#1{\@startsection {subsection}{2}{0pt}%
{-3.5ex \@plus -1ex \@minus -.2ex}%
{1ex \@plus.2ex}%
{\bf\mathversion{bold}}{#1}}

\def\subsubsection#1{\@startsection{subsubsection}{3}{0pt}%
{\medskipamount}%
{-10pt}%
{\normalsize\itshape}{\kern-2.2ex. #1.}}


\def\blfootnote{\xdef\@thefnmark{}\@footnotetext}

\makeatother

\def\note#1{\par\smallskip%
\noindent%
\llap{$\boldsymbol\Longrightarrow$}%
\fbox{\vtop{\hsize=0.98\hsize\parindent=0cm\small\rm #1}}%
\rlap{$\boldsymbol\Longleftarrow$}%
\par\smallskip}

\renewcommand{\cite}{\citet}

\def\^#1{\ifmmode {\mathaccent"705E #1} \else {\accent94 #1} \fi}
\def\~#1{\ifmmode {\mathaccent"707E #1} \else {\accent"7E #1} \fi}

\edef\-#1{\noexpand\ifmmode {\noexpand\bar{#1}} \noexpand\else \-#1\noexpand\fi}
\def\>#1{\vec{#1}}
\def\.#1{\dot{#1}}
\def\wh#1{\widehat{#1}}
\def\wt#1{\widetilde{#1}}
\def\atop{\@@atop}

\renewcommand{\leq}{\leqslant}
\renewcommand{\geq}{\geqslant}
\renewcommand{\phi}{\varphi}
\newcommand{\eps}{\varepsilon}
\newcommand{\D}{\Delta}
\newcommand{\eq}{\eqref}

\newcommand{\dtv}{\mathop{d_{\mathrm{TV}}}}
\newcommand{\dw}{\mathop{d_{\mathrm{W}}}}
\newcommand{\dk}{\mathop{d_{\mathrm{K}}}}

\newcommand{\bigo}{\mathrm{O}}
\newcommand{\lito}{\mathrm{o}}
\newcommand{\hence}{\Rightarrow}
\newcommand{\toinf}{\to\infty}
\newcommand{\tozero}{\to0}
\newcommand{\longto}{\longrightarrow}
\newcommand{\Var}{\mathop{\mathrm{Var}}\nolimits}
\newcommand{\law}{\mathscr{L}}

\newcommand{\ER}{Erd\H{o}s-R\'{e}nyi}


\newcommand{\SL}{\mathrm{SL}}
\newcommand{\Exp}{\mathrm{Exp}}
\newcommand{\Ber}{\mathrm{Ber}}
\newcommand{\Geo}{\mathrm{Geo}}
\newcommand{\Beta}{\mathrm{Beta}}
\def\equ#1{^{(#1)}}
\def\zer#1{^{[#1]}}
\def\dsb#1{^{\langle#1\rangle}}
\newcommand{\eqd}{\stackrel{d}{=}}
\newcommand{\im}{\mathrm{i}}
\newcommand{\convd}{\stackrel{d}{\longrightarrow}}
\newcommand{\bu}{\mathbf{u}}
\newcommand{\bx}{\mathbf{x}}
\newcommand{\bzero}{\mathbf{0}}
\newcommand{\by}{\mathbf{y}}
\newcommand{\bw}{\mathbf{w}}
\newcommand{\te}{\mathbf{e}}
\def\sp#1{^{(#1)}}
\newcommand{\MVN}{\mathrm{MVN}}


\def\BC{\mathrm{BC}}

%%%%% My own commands and packages%%%%%%%%%%%%%%%%%%%

\usepackage{nomencl}
\makenomenclature

\usepackage{mwe}
\usepackage{enumitem}
\usepackage{xcolor}
\usepackage{mathtools} % for \bigtimes 


\newcommand{\R}{\mathbb{R}}
\newcommand{\N}{\mathbb{N}}
\newcommand{\E}{\mathbb{E}}
\newcommand{\Hn}{\mathbb{H}_{1,+}} % for nonnegative harmonics with h(x0) = 1

\renewcommand{\Var}{\operatorname{Var}}
\newcommand{\Cov}{\operatorname{Cov}}

\renewcommand{\P}{\mathbb{P}}
\newcommand{\bt}{\mathbf{t}}
\newcommand{\bs}{\mathbf{s}}
\newcommand{\Xn}{(X_n)_{n \in \mathbb{N}}}

\newcommand{\induced}{\text{ind}}

\newcommand{\red}{\textcolor{red}}
\newcommand{\pink}{\textcolor{pink}}

\DeclareMathOperator*{\argmin}{arg\,min}


% indicator function
\newcommand{\ind}{\mathbf{1}}

% tilde commands
\newcommand{\tx}{\tilde{x}}
\newcommand{\tX}{\tilde{X}}
\newcommand{\ty}{\tilde{y}}
\newcommand{\tK}{\tilde{K}}
\newcommand{\tE}{\tilde{E}}

% Independence
\newcommand{\indep}{\perp\!\!\!\perp}

\newcommand{\zack}[1]{\textcolor{Violet}{Zack comment: #1}}
\newcommand{\daniel}[1]{\textcolor{Emerald}{Daniel comment: #1}}

\newtheorem*{theorem*}{Theorem}
\newtheorem*{lemma*}{Lemma}